\chapter{The Geometry of Consciousness}
\label{ch:geometry}

\papernote{Source paper: \emph{On the Geometry of Consciousness} --- \texttt{docs/geometry/}}

\begin{chapterabstract}
Consciousness cannot be directly measured: any attempt to access conscious content transforms it through the observer's own consciousness. This is the observer-participancy problem---more fundamental than Heisenberg uncertainty, establishing privacy as an intrinsic property of consciousness rather than a limitation of measurement technology. However, consciousness possesses measurable \emph{geometry}: the invariant mathematical structure at the confluence of perception flux and thought geometry. The Privacy Axiom is stated and proved. The Confluence Theorem establishes the ``NOW'' as the intersection point $\NOW = f(\tau_p, \tau_t)$ of perception and thought decay curves. Quantitative clinical metrics (PLV, $\mathcal{C}_{\mathrm{confluence}}$, $\mathcal{S}_{\mathrm{equilibrium}}$) distinguish healthy consciousness, dissociative states, and disorders of consciousness.
\end{chapterabstract}

\input{../docs/geometry/geometry-of-consciousness.tex}
